\section{iPRA adequacy and a clear route to complete Lean formalization}
\label{sec:formalization}

\subsection{The internal statements}

Write $\operatorname{HurCert}(B,p,q,s)$ for the conjunction of the
positive-denominator, coprimality, threshold and strict rational tests in
\eqref{eq:hw-certificate}.  After rational arithmetic has been replaced
by its integer-code operations, this is a quantifier-free predicate.
For each fixed pair of primitive recursive definitions $x_s,I$, the
argument constructs primitive recursive output terms $p(B),q(B),s(B)$.
Its arithmetic conclusion is the iPRA sequent
\[
 \operatorname{Name}(x)\wedge\operatorname{Apart}(x,I)
 \ \longrightarrow\
 \forall B\;\operatorname{HurCert}(B,p(B),q(B),s(B)).
\]
Here $\operatorname{Name}$ and $\operatorname{Apart}$ are ordinary
first-order universal formulas about the fixed rational-name functions;
the symbols $x,I$ in this display name those definitions and are not
variables of a new function sort.  The monotonic running maximum for $I$
is part of the primitive recursive construction.

For each fixed rational $\psi,D$, let
$\operatorname{Mon}(\psi)$ express nonnegative monotonic values and let
$\operatorname{Div}(\psi,D)$ be (K-D).  The metric construction supplies
primitive recursive terms $U(N,m),M(N,m)$ and the sequent
\[
 \operatorname{Mon}(\psi)\wedge\operatorname{Div}(\psi,D)
 \ \longrightarrow\
 \forall N,m\,[N\ge1\to\operatorname{Cover}(N,m,U(N,m),M(N,m))].
\]
The predicate $\operatorname{Cover}$ tests that the interval list is
correctly labelled by reduced fractions, every label has $N<q\le M$,
each radius fits the target radius, and its rational residual length is
less than $2^{-m}$.  Each finite-list universal condition is a Boolean
fold, so the final predicate is quantifier-free.  The convergence
sequent similarly asserts the finite length bound under (K-C).

One can first state these as unconditional finite-transcript theorems:
all name evaluations, modulus evaluations and algebraic checks used by a
run are supplied as a finite table; if the table passes the requisite
tests, the output passes the stated certificate test.  Instantiating
the table entries with fixed primitive recursive definitions gives the
schemata above.  This separates arithmetic correctness from any
assumption about the source of the data.

Every code operation must be total on natural numbers.  On malformed
fractions or lists, use a designated default value and a validity bit.
The proofs show that such defaults are never used under the hypotheses.
The reference Python code raises exceptions on invalid inputs for
readability; a literal compilation into PR syntax instead uses these
totalized operations.

\subsection{The induction audit}

The following table gives the finite invariants needed to implement the
human proof.  All comparisons are integer comparisons after clearing
known positive denominators.  Bounded quantifiers in the middle column
are to be compiled into Boolean folds as in Section~\ref{sec:foundations}.

\begin{longtable}{p{0.20\textwidth}p{0.47\textwidth}p{0.23\textwidth}}
\toprule
Computation & Invariant & Induction index\\
\midrule
\endhead
Integer and fraction routines & Sign/denominator validity; cross-product
identities; bounded division remainder and gcd witnesses & Explicit
arithmetic loop counter\\
Convergents & Determinant $\pm1$; positive denominators; growth; stored
rational enclosure has the indicated finite prefix & Prefix length\\
Digit extraction & Fractional-linear denominator bounded away from zero;
image enclosure avoids every relevant integer; endpoint floors agree &
Prefix length and bounded comparison list\\
Hurwitz selection & Chosen pair has $5C^2<1$; precision has its prescribed
margin; selected rational error satisfies the weighted bound & Finite
candidate scan; algebraic identities\\
Counting rational centers & The accumulated count equals the number of
accepted pairs; common-divisor indicators dominate the bad-pair bit &
Bounds on $p,q,d$\\
Dyadic estimates & Running harmonic block sum and telescoping
inverse-square sum have the displayed bounds; $2^k\ge k^2$ for $k\ge4$ &
Integer bound and dyadic exponent\\
Endpoint partition & Sorted endpoints preserve the multiset; each cell
has the correct coverage and incidence bits & Sort and finite cell index\\
Second moment & Accumulated first and second moments satisfy the finite
incidence identities; sum of squares is nonnegative & Cell and pair indices\\
Mass block & Finite grouping of the original tail into dyadic blocks;
capped partial sum dominates the capped total & Dyadic block index\\
Gap filling & Old gaps are processed once; added intervals lie inside
their parent gap and have correct labels; total residual length contracts &
Frozen gap list length, then fixed round count\\
Null covers & Padded list length obeys its geometric budget; finite
rational locator chooses a good interval or a forbidden padded interval &
Finite list and cover enumeration index\\
\bottomrule
\end{longtable}

This table is not a claim that an arbitrary algorithm with a bounded
loop automatically has an iPRA proof.  The algebraic identities,
indicator inequalities, and preservation arguments proving these
particular invariants are given in the preceding sections.  Their
quantifier-free forms are precisely the instances of induction needed.
The finite sums, products and list transformations are primitive
recursive by their displayed recurrences and bounded enumerations.
The large number $16384(m+1)$ of outer rounds is prescribed before the
run.  A round is a fixed primitive recursive transformation of a finite
state; iterating it that many times is again primitive recursion.
New gaps are not processed by an uncontrolled recursive call in the
same round.

The real-number shorthand in the Hurwitz proof has a finite rational
interface, as shown at the end of its extraction argument.  A formal
implementation should prove that interface first and use name coherence
only to supply its enclosures.  It should not build a general classical
real-number theory inside the object language.

\subsection{What was verified in the repository}

The inspected repository snapshot is
\begin{center}
\texttt{34c4c35edc8e648e57587c6faf3a5cfdd7879777}.
\end{center}
Its toolchain is Lean \texttt{v4.32.0}; the mathlib manifest resolves
\texttt{v4.32.0} to the following revision:
\begin{center}
\texttt{81a5d257c8e410db227a6665ed08f64fea08e997}.
\end{center}
The relevant namespace is \texttt{LeanProofs.ConciseIPRA}.  Its source
directory is
\begin{quote}\small
\path{Logic/PrimitiveRecursiveArithmetic/ConciseIPRA/Lean/ConciseIPRA/}.
\end{quote}
The source inspection and exact pinned URLs are recorded in
\path{sources/repository_audit.txt} and \path{sources/source_manifest.json}.

The useful existing interfaces are concrete:\begin{itemize}
\item \texttt{PR.lean} defines the primitive recursive syntax
\texttt{PR}, term syntax \texttt{PTm}, axiom family \texttt{PRAAx}, and
the actual object-theory derived rule \texttt{PRA.qfInd}.
\item \texttt{TeqReflect.lean} supplies characteristic terms
\texttt{chi}, the reflection result \texttt{chi\_iff}, and
\texttt{Teq.qfInd}.  Its quantified formula cases have dummy
characteristic values, so it must not be applied directly to bounded
quantified formulas.  The bounded-fold compiler described above is needed.
\item \texttt{TeqPair.lean} and \texttt{TeqRec.lean} provide pairing and
primitive-recursion machinery, including \texttt{Tm.prc},
\texttt{Q.prc\_zero} and \texttt{Q.prc\_succ}.
\item \texttt{BiInterpretation.lean} proves actual proof translations
between \texttt{TeqAx} and \texttt{PRAAx}.
\item \texttt{NativeBiInterpretation.lean} provides
\texttt{pra\_iff\_native}, translating a \texttt{PRAAx} derivation to
the concise native presentation.
\end{itemize}

The last bridge has the following source-level shape, with its printed
notation expanded informally:
\begin{lstlisting}
(Gamma proves phi in PRAAx) <->
  NPrf (Gamma.map NFm.ofPRA) (NFm.ofPRA phi)
\end{lstlisting}
It can be applied after an actual object-theory proof has been built;
it does not turn an arbitrary ambient Lean theorem into such a proof.
The repository README reports completed foundations without additional
axioms or \texttt{sorry}, with metatheoretic dependencies including
\texttt{propext}, \texttt{Quot.sound}, and, in part of the native chain,
\texttt{Classical.choice}.  Those are dependencies of the Lean
metatheory; the embedded iPRA derivations still use their explicitly
defined intuitionistic rules.  We inspected the relevant source and did
not replay repository builds in this environment.

The \path{docs/drafts/} listing contains a Dirichlet/Kronecker archive,
but its binary contents were not inspected, so no theorem from it is
assumed.  An inspected ordinary Lean finite simultaneous-Dirichlet
module gives a useful pigeonhole pattern; it uses a classical block and
is not itself an embedded iPRA derivation.  Neither main proof in this
report depends on that module.

\subsection{Formalization milestones and completion criteria}

\paragraph{Milestone 1: executable finite arithmetic.}
Implement signed integer and fraction codes, finite list operations,
bounded search, gcd, sorting and interval content.  Prove the primitive
recursive closure certificates as well as their semantic correctness.
The Python artifact gives an executable reference for endpoint behavior
and certificate values; the Lean code should have independent proofs.

\paragraph{Milestone 2: reference Lean theorems.}
Prove the denominator polynomial lemma and weighted-error selector over
$\Z$ and $\Q$.  Prove the finite center-count, overlap and second-moment
lemmas using finite sets or lists and exact rational sums.  Polynomial
normalization tactics can discharge the displayed identities; all
inequalities must be proved from the explicit sign hypotheses.  Build
the rational-enclosure version of the continued-fraction extractor.
Do not use an imported real Hurwitz or metric approximation theorem to
replace these finite lemmas.

\paragraph{Milestone 3: the constructive algorithms.}
Formalize the fixed-round metric state transformation, its local
coverage invariant and contraction.  Formalize the name/modulus
interfaces and the padded-cover locator.  At this point an ambient Lean
theorem proves each finite specification.  This is a substantial useful
checkpoint, but is still separate from the iPRA certification.

\paragraph{Milestone 4: actual iPRA derivations.}
Compile each fixed arithmetic definition to \texttt{PR}/\texttt{PTm}.
Prove bounded-fold reflection inside \texttt{PRAAx}; then implement
the invariants in the table using \texttt{PRA.qfInd}, substitution,
equality and intuitionistic natural deduction.  The final deliverable
for this milestone is a Lean term whose type is an actual
\texttt{Prf PRAAx} sequent for the two displayed schemata (or a
metatheorem generating that term from the fixed primitive recursive
definitions and their hypotheses).  An abstract record assuming these
derivations would only name the remaining task.

\paragraph{Milestone 5: concise native presentation.}
Apply \texttt{pra\_iff\_native} to the constructed derivation.  This uses
the already implemented bridge and introduces no additional number
theory.  Optional classical-comparison modules may interpret the finite
bounds as Lebesgue statements over $\mathbb R$; keep their dependencies
separate from the iPRA result.

The included \path{lean/Blueprint.lean} records finite interfaces and
explicit obligations.  \path{lean/FORMALIZATION_PLAN.txt} describes the
implementation order and repository integration points.  These files
do not claim completion of any of the above milestones beyond their
stated definitions.  On a machine with the pinned repository environment,
the appropriate first check is
\begin{lstlisting}
lake env lean /absolute/path/to/lean/Blueprint.lean
\end{lstlisting}
followed by the separate module builds as the obligations are discharged.

\subsection{Executable checks supplied with this report}

The command \texttt{python3 code/verify.py} regenerates the exact
\path{code/verification.json} record.  Its checks include 9,830 denominator
triples; six Hurwitz certificates for the golden ratio and $\sqrt2$;
1,000 seeded union comparisons against an independent endpoint-cell
calculation; 1,000 intersection comparisons; 1,000 cap inequalities;
90 local overlap estimates; and explicit endpoint edge cases.  All
calculations use integers and exact fractions.  These computations
check implementations and examples; they do not replace the general
proofs.

For instance the golden-ratio example with threshold $B=10$ returns
$p/q=233/144$ and a precision-$39$ rational certificate whose left side is
\[
 \frac{18446391806604141645}{18446744073709551616}<1.
\]
This tests a case close to the sharp constant.  The finite reduced union
for $\psi(q)=1/q$ and $5\le q\le80$ has 1,960 labelled intervals and
exact covered length
\[
 \frac{55336826409}{59204622400}.
\]
Neither this sample denominator cutoff nor the moderate local tests are
claimed to instantiate the astronomical worst-case cutoff in the proof.
The full general residual-gap algorithm and the general modulus-based
continued-fraction extractor are specified in the manuscript; they have
not been executed end to end by the sample program.
